\documentclass[11pt,a4paper]{article}

\usepackage[T1]{fontenc}
\usepackage[utf8]{inputenc}
\usepackage{amsmath,amssymb}
\usepackage[a4paper,margin=3cm]{geometry}
\usepackage[colorlinks=true,linkcolor=black,citecolor=black,urlcolor=blue!55!black]{hyperref}
\usepackage{setspace}

\pagestyle{empty}

\begin{document}

\begin{center}
{\large\bfseries University of Parma}\\[2pt]
{\normalsize Department of Mathematical, Physical and Computer Sciences}\\[4pt]
{\normalsize Bachelor's Degree in Physics}\\[18pt]

{\Large\bfseries Quantum simulation of molecular spin systems\\ on quantum hardware}\\[18pt]

{\normalsize Candidate: \textbf{Samuele Schiavi}}\\[4pt]
{\normalsize Supervisor: \textbf{Prof. Alessandro Chiesa}}\\[18pt]
{\normalsize\bfseries Abstract}
\end{center}

\vspace{6pt}
\onehalfspacing
\noindent
Molecular magnets are privileged platforms for studying collective quantum
phenomena, yet their classical simulation quickly becomes intractable as the
number of spins increases. Quantum computers, by mapping each spin-$1/2$
directly onto a qubit, offer a natural route around this bottleneck. This
work presents the quantum simulation, on both ideal and realistically noisy
simulators, of spin-$1/2$ systems coupled through a Heisenberg interaction
with an anisotropic Dzyaloshinskii--Moriya term, starting from the simplest
case (the dimer) up to a frustrated three-spin ring. The ground state is
determined with the Variational Quantum Eigensolver, comparing a generic
hardware-efficient ansatz with ans\"atze built on the physical symmetries of
the problem; time evolution is implemented through a Suzuki--Trotter
decomposition, and dynamical spin-spin correlations are extracted with a
Hadamard-test circuit. The full pipeline is then repeated under a noise
model calibrated on real quantum hardware, identifying the number of
Trotter steps that balances the approximation error against the error
accumulated from noise, and testing its stability against noise-aware
reoptimization and asymmetric readout error. The results show that an
ansatz aligned with the symmetries of the Hamiltonian drastically reduces
the number of parameters required with respect to a generic ansatz, an
advantage that persists, at increasing cost, when moving from two to three
spins --- a design criterion that is a promising candidate for extension to
larger molecular systems.

\vspace{12pt}
\noindent\textbf{Keywords:} quantum simulation; spin systems; Variational
Quantum Eigensolver (VQE); Suzuki--Trotter decomposition; quantum noise;
IBM Quantum hardware.

\end{document}
